\(\mathsf{DCTFTR}\) may return the fast-path identifying functional \(F_B(\mathcal P)/F_E(\mathcal P)\in\mathscr R(\mathcal P)\) only with a certificate that \(F_E>0\) throughout \(\mathcal C(\mathcal P)\).  Otherwise \(\mathsf{DCTFTR\text{-}CAN}\) uses \(\Lambda=\Theta_G^{\mathrm{can}}(\mathcal P)\).  For an unconditional target cell, its exact output is
\[
\left[
\min_{x\in\Lambda}\Phi_{T,\mathrm{can}}^\star(x^\star),
\max_{x\in\Lambda}\Phi_{T,\mathrm{can}}^\star(x^\star)
\right].
\]
For a conditional target, let
\[
e_-=\min_{x\in\Lambda}\Phi_{E,\mathrm{can}}^\star(x^\star),
\qquad
e_+=\max_{x\in\Lambda}\Phi_{E,\mathrm{can}}^\star(x^\star).
\]
The output is zero evidence when \(e_+=0\), is not uniformly defined when \(e_-=0<e_+\), and only when \(e_->0\) is the exact conditional interval
\[
\left[
\min_{x\in\Lambda}
\frac{\Phi_{B_\rho\cap E,\mathrm{can}}^\star(x^\star)}
{\Phi_{E,\mathrm{can}}^\star(x^\star)},
\max_{x\in\Lambda}
\frac{\Phi_{B_\rho\cap E,\mathrm{can}}^\star(x^\star)}
{\Phi_{E,\mathrm{can}}^\star(x^\star)}
\right].
\]
Equal endpoints give an exact algebraic identification certificate and unequal endpoints give canonical original-alphabet paired-model witnesses.  Separately, local district-linear conclusions use
\[
R_D(m_D)=\operatorname{row}(A_D)+\operatorname{span}\{e_\omega:\omega\in Z_D(m_D)\}
=\operatorname{span}\{q-q':q,q'\in\mathcal K_D(m_D)\}^{\perp},
\]
with sharp focal realization only under \(\mathsf{BASep}_D(\mathcal P)\) or its stated stronger sufficient checks; cross-district global cells remain in the canonical semialgebraic program.